% ============================================================================
% Fichier  : partie0/P00_nuit_bafoussam.tex
% Titre    : La Nuit du 14 Mars à Bafoussam
% Auteur   : MINKA MI NGUIDJOÏ Thierry Emmanuel
% Version  : 1.0 -- novembre 2024
% Corrigés : corriges/Corr_P00.tex
% ============================================================================
% Faits posés (bible des cas, Claude outputs/01_fiche_plan_bible.md) :
% NYANGA CRÉDIT, Bafoussam, 22 agences, 210 000 clients ; 23 virements
% frauduleux de 02 h 05 à 03 h 31 le 14/03/2025, 47 300 000 FCFA, 11
% portefeuilles ; exfiltration de 3,2 Go de SRV-FIC01 vers 203.0.113.47 de
% 01 h 12 à 03 h 47 ; faille VULN-2024-4471 (fictive) de la passerelle VPN,
% correctif du 15/12/2024, KEV le 19/12/2024, premier accès le 02/03/2025.
% ============================================================================

\chapter{La Nuit du 14 Mars à Bafoussam}
\label{chap:prologue}
\chaptermark{La nuit du 14 mars}

\epigraph{%
    « Ce que l'on fait dans la première heure\\
    décide de ce que l'on pourra prouver dans un an. »%
}{Principe du cours}

\niveauChapitre{Introductif}{%
    Aucun prérequis~: ce prologue ouvre le cours et en pose les questions.
}

\begin{boxobjectifs}
\begin{itemize}
    \item Reconnaître, dans un incident réel, les questions qui relèvent de
          l'investigation numérique et celles qui relèvent de l'analyse de
          vulnérabilités.
    \item Distinguer les gestes qui préservent la preuve de ceux qui la
          détruisent pendant la première heure.
    \item Séparer, dans un récit d'incident, la constatation de l'analyse.
    \item Coter une conclusion d'enquête \OUI{}, \NON{} ou \NEK{}.
    \item Situer les métiers qui interviennent, de l'alerte à la correction.
\end{itemize}
\end{boxobjectifs}

\bigskip

% ============================================================================
\section*{Avant de commencer : un vendredi matin}
% ============================================================================

Bafoussam, vendredi 14 mars 2025, 7~h~30. Au service des paiements de NYANGA
CRÉDIT\footnote{NYANGA CRÉDIT, SOLITRANS LOGISTIQUE, leurs dirigeants, leurs
prestataires, les adresses, les dates et les montants sont fictifs~;
l'identifiant VULN-2024-4471 est fictif. Les mécanismes techniques, les
outils, les normes et les textes cités sont exacts.}, une agente de
permanence ouvre le rapprochement de la nuit. Elle compte vingt-trois
virements sortants passés entre 2~h~05 et 3~h~31, tous par l'interface qui
relie l'établissement aux opérateurs de paiement mobile, tous vers onze
portefeuilles que personne ne connaît. Le total atteint 47\,300\,000 francs
CFA. Aucun agent n'était en poste à ces heures-là.

NYANGA CRÉDIT est un établissement de microfinance de deuxième catégorie,
installé à Bafoussam, présent dans vingt-deux agences de l'Ouest et du
Littoral, au service d'environ 210\,000 clients. Son informatique tient dans
une salle du siège~: un contrôleur de domaine, un serveur comptable, un
serveur de fichiers, l'interface de paiement mobile, et une passerelle
d'accès à distance qui permet aux agences et aux techniciens de se connecter
au siège.

À 8~h~15, le directeur des systèmes d'information, M.~KAMGA, arrive. Le
serveur de fichiers ralentit depuis la veille. Son premier réflexe est de le
redémarrer, puis de restaurer la sauvegarde de la veille. La responsable de la
sécurité, Mme~NDOUMBE, l'arrête~: elle fait débrancher le câble réseau du
serveur, sans l'éteindre, et demande que personne ne touche plus aux
machines du siège avant son retour. Elle revient vingt minutes plus tard avec
une clé USB qui contient un outil de capture de la mémoire vive.

\medskip

Ce prologue raconte les premières heures de l'incident et les questions
qu'elles posent. Les seize chapitres qui suivent y répondent une à une.

\begin{boxalert}
\textbf{Violation typique.} Redémarrer, nettoyer ou restaurer une machine
compromise avant d'avoir préservé ce qu'elle contient.

\smallskip
\textbf{Règle~:} on isole avant d'agir, on préserve avant d'analyser, et on
consigne chaque geste avec son heure et son auteur. Une machine redémarrée a
perdu sa mémoire vive, ses connexions ouvertes et une partie de ses
journaux~; une machine restaurée a perdu tout le reste.
\end{boxalert}

\separateur

% ============================================================================
\section{Ce que l'on sait à 9 heures}
\label{sec:p00-9h}
% ============================================================================

À 9~h~10, le directeur général, M.~TCHOUA, réunit une cellule de crise~:
M.~KAMGA, Mme~NDOUMBE et la juriste de l'établissement, Me~ETOUNDI. Mme~NDOUMBE
pose sur la table ce qui est établi, et seulement cela.

\begin{tableaucours}{L{5.2cm}L{4.4cm}c}{%
    \tableauHeader{Constatation} &
    \tableauHeader{Source} &
    \tableauHeader{Cote}
}
23 virements entre 2~h~05 et 3~h~31, 47,3~M~FCFA, 11 portefeuilles &
Journal de l'interface de paiement & \OUI{} \\
Les virements ont utilisé la clé d'accès de l'établissement à l'interface &
Journal de l'interface de paiement & \OUI{} \\
Aucun agent n'était connecté au siège entre 1~h et 5~h & Registre des
présences, journal de la passerelle & \OUI{} \\
Le serveur de fichiers a émis un volume sortant inhabituel dans la nuit &
Supervision du lien Internet & \OUI{} \\
Ce volume correspond à des fichiers de clients & Aucune pièce à ce stade &
\NEK{} \\
L'attaquant est entré par la passerelle d'accès à distance & Hypothèse du
DSI & \NEK{} \\
L'attaque vient d'un employé & Rumeur dans les couloirs & \NEK{} \\
\end{tableaucours}
\vspace{-0.8em}
\begin{center}{\small\itshape Tableau~0.1. Le point de 9~h~10}\end{center}

Trois lignes sont cotées \NEK{}~: aucune preuve ne permet encore de les
établir ou de les écarter. Une hypothèse, même plausible, ne devient pas une
constatation parce qu'un responsable la formule avec assurance. Mme~NDOUMBE
le dit à la cellule~: « Nous savons ce qui est sorti du compte. Nous ne
savons pas encore par où ils sont entrés. »

\begin{boxdef}
\textbf{\termecle{Constatation}}~: fait vérifiable, rattaché à une source
identifiée (un journal, une image, une capture), que toute personne
compétente retrouverait en refaisant l'examen.

\smallskip
\textbf{\termecle{Analyse}}~: interprétation d'une ou de plusieurs
constatations~; elle porte un niveau de confiance (faible, moyenne, élevée)
et peut être révisée par une constatation nouvelle.
\end{boxdef}

\begin{boxdef}
\textbf{\termecle{Cote OUI, NON, NEK}}~: réponse à une question d'enquête
fermée. \OUI{}~: la proposition est établie par une preuve. \NON{}~: elle est
contredite par une preuve. \NEK{} (\textit{not enough knowledge})~: les
preuves disponibles ne permettent pas de trancher. L'instrument vient du
programme de recherche de l'auditologie~\cite{minka2026nek}.
\end{boxdef}

% ============================================================================
\section{Les gestes de la première heure}
\label{sec:p00-gestes}
% ============================================================================

Ce que Mme~NDOUMBE fait entre 8~h~15 et 10~h~30 n'a rien d'improvisé~: elle
applique un ordre que l'on retrouve dans la RFC~3227~\cite{rfc3227} et dans
la norme ISO/IEC~27037~\cite{iso27037_2012}, celui de la volatilité. On
recueille d'abord ce qui disparaît le plus vite.

\begin{tableaucours}{cL{3.6cm}L{4.0cm}L{3.4cm}}{%
    \tableauHeader{Heure} &
    \tableauHeader{Geste} &
    \tableauHeader{Ce qu'il préserve} &
    \tableauHeader{Ce qu'il évite}
}
8~h~20 & Débrancher le câble réseau du serveur de fichiers, sans l'éteindre &
La mémoire vive, les processus, les connexions récentes & La poursuite de
l'exfiltration, l'effacement à distance \\
8~h~25 & Interdire toute connexion aux serveurs du siège & L'état des
disques et des journaux & Les modifications par l'équipe elle-même \\
8~h~40 & Désactiver la passerelle d'accès à distance & Les journaux de la
passerelle, figés & Un nouvel accès \\
10~h~02 & Capturer la mémoire du contrôleur de domaine (16~Go), calculer
son empreinte & Les processus et les identifiants en mémoire & La perte au
prochain redémarrage \\
10~h~30 & Ouvrir un registre des actions~: heure, auteur, geste, résultat &
La traçabilité de chaque intervention & La contestation ultérieure \\
\end{tableaucours}
\vspace{-0.8em}
\begin{center}{\small\itshape Tableau~0.2. Les gestes de la première heure}\end{center}

Chaque geste a un coût. Débrancher le serveur arrête aussi le partage de
fichiers des agences~; désactiver la passerelle prive les techniciens
d'accès à distance. Le DSI le fait remarquer. La réponse de Mme~NDOUMBE tient
en une phrase~: une heure d'interruption se rattrape, une preuve détruite
ne se reconstitue pas.

\begin{boxobjet}{Capture mémoire de SRV-AD01, 14 mars 2025, 10~h~02}
\textbf{Source~:} contrôleur de domaine SRV-AD01, salle serveurs du siège,
Bafoussam.

\textbf{Pièce~:} fichier \texttt{SRV-AD01\_20250314\_1002.raw}, 16~Go,
copié sur un disque externe neuf, étiqueté et placé sous enveloppe.

\textbf{Empreinte~:} SHA-256 calculée à 10~h~41 et inscrite au registre,
contresignée par M.~KAMGA.

\textbf{Ce qu'elle peut prouver~:} les processus actifs à 10~h~02, les
connexions réseau ouvertes, les traces d'outils exécutés en mémoire.

\textbf{Ce qu'elle ne prouve pas~:} ce qui s'est passé avant le dernier
redémarrage de la machine, ni l'identité d'une personne.
\end{boxobjet}

% ============================================================================
\section{L'enquête : ce qu'elle établira}
\label{sec:p00-enquete}
% ============================================================================

L'après-midi même, NYANGA CRÉDIT mandate le cabinet SANKOFA EXPERTISE, à
Yaoundé, pour conduire la réponse à l'incident~; Mme~MBARGA dirige la
mission. Le lendemain, l'établissement dépose plainte au parquet de
Bafoussam. Le 17~mars, le procureur de la République réquisitionne un expert
indépendant, M.~ABOUBAKAR, expert près la cour d'appel de l'Ouest. Deux
enquêtes avancent désormais en parallèle~: celle du cabinet, au service de
l'établissement, et celle de l'expert, au service de la justice. Elles
examinent les mêmes pièces et n'ont pas les mêmes obligations.

Trois semaines plus tard, les deux enquêtes établissent la même séquence.

\begin{boxscenario}{Opération NYANGA. Ce que l'enquête établit}
\textbf{Constatations.}
\begin{itemize}
    \item Le 2~mars à 23~h~48, le journal de la passerelle enregistre une
          première session du compte de service \texttt{svc\_backup}, depuis
          l'adresse 203.0.113.18, sans échec d'authentification préalable.
    \item Le 9~mars à 1~h~37, un fichier \texttt{lsass.dmp}, copie de la
          mémoire du processus qui gère les authentifications Windows, est
          créé dans un dossier temporaire de SRV-AD01~; la capture mémoire
          du 14~mars montre encore un outil d'accès à distance actif et une
          connexion vers 203.0.113.47.
    \item Le 12~mars, le journal de SRV-CPT01 montre la lecture du fichier
          de configuration qui contient la clé d'accès à l'interface de
          paiement.
    \item Le 14~mars, de 1~h~12 à 3~h~47, la capture réseau montre 3,2~Go
          envoyés de SRV-FIC01 vers 203.0.113.47.
    \item La passerelle fonctionnait avec une version de micrologiciel
          antérieure au correctif publié le 15~décembre 2024 pour la faille
          VULN-2024-4471, qui permet de contourner l'authentification.
\end{itemize}

\textbf{Analyse} (confiance élevée). L'attaquant est entré par la
passerelle, en exploitant VULN-2024-4471 pour ouvrir une session sous un
compte de service~; il a récupéré des identifiants d'administration en
mémoire, lu la clé de l'interface de paiement, copié des fichiers, puis
déclenché les virements.

\textbf{Ce qui reste \NEK{}.} L'identité des auteurs. Les adresses
203.0.113.18 et 203.0.113.47 désignent des machines, pas des personnes.
\end{boxscenario}

La dernière constatation change la nature de l'affaire. Le correctif était
disponible depuis 89~jours le jour de l'attaque, et depuis 77~jours lors du
premier accès~; le 19~décembre, la faille avait été inscrite au catalogue
des vulnérabilités exploitées tenu par l'agence américaine CISA. L'enquête,
partie d'un virement frauduleux, aboutit à une faiblesse connue, publiée et
corrigeable. Pourquoi n'a-t-elle pas été corrigée~? Parce que la passerelle
ne figurait pas dans l'inventaire de l'outil de balayage de vulnérabilités~:
personne ne l'a jamais analysée.

% ============================================================================
\section{Deux disciplines, une même affaire}
\label{sec:p00-deux}
% ============================================================================

L'affaire NYANGA fait se rejoindre deux disciplines que l'on enseigne
souvent séparément.

\begin{tableaucours}{L{3.2cm}L{4.6cm}L{4.6cm}}{%
    \tableauHeader{} &
    \tableauHeader{Investigation numérique} &
    \tableauHeader{Analyse de vulnérabilités}
}
Point de départ & Une trace, un incident & Une faiblesse connue, un actif \\
Question & Que s'est-il passé, comment, depuis quand~? & Où sommes-nous
exposés, et que corriger d'abord~? \\
Matière & Images de disque, mémoire, captures, journaux & Inventaires,
résultats de balayage, bulletins, scores \\
Produit & Un rapport qui peut servir de preuve & Un plan de correction
suivi \\
Temps & Après l'incident & Avant l'incident, en continu \\
\end{tableaucours}
\vspace{-0.8em}
\begin{center}{\small\itshape Tableau~0.3. Deux disciplines}\end{center}

Elles se nourrissent l'une l'autre. L'enquête révèle les faiblesses que
l'attaquant a réellement empruntées~; elle dit au programme de gestion des
vulnérabilités ce qu'il a manqué. Le programme, en retour, réduit le nombre
d'enquêtes à conduire et fournit aux enquêteurs l'inventaire et l'historique
des versions qui leur permettent de dater une exposition. À Douala, un autre
acteur du cours, SOLITRANS LOGISTIQUE, montrera dans la partie~IV comment on
construit un tel programme sur 1\,380 actifs.

\begin{boxinfo}
\textbf{Ce que le cours ne traite pas.} L'exploitation des failles relève
de l'UE MSN-2S3-SEC03 (Ethical Hacking et Pentesting)~: ce cours apprend à
identifier, coter et faire corriger une vulnérabilité, pas à s'en servir.
La protection des données personnelles exposées par l'exfiltration relève de
l'UE MSN-2S3-JUR01. L'audit et la gestion du risque relèvent de l'UE
MSN-2S3-SEC08.
\end{boxinfo}

% ============================================================================
\section{Les questions du cours}
\label{sec:p00-questions}
% ============================================================================

Chaque chapitre répond à l'une des questions que la cellule de crise s'est
posées, ou aurait dû se poser.

\begin{tableaucours}{L{7.4cm}L{4.8cm}}{%
    \tableauHeader{Question de la cellule} &
    \tableauHeader{Chapitre}
}
Qu'est-ce qu'enquêter sur des traces numériques~? &
Chapitre~\ref{chap:investigation} \\
Comment faire pour que ce que l'on trouve serve de preuve~? &
Chapitre~\ref{chap:preuve} \\
Comment copier un disque ou une mémoire sans les altérer~? &
Chapitre~\ref{chap:acquisition} \\
Que contient le disque du serveur comptable~? & Chapitre~\ref{chap:disque} \\
Que révèle la mémoire du contrôleur de domaine~? &
Chapitre~\ref{chap:memoire} \\
Qu'est-il sorti, et vers où~? & Chapitre~\ref{chap:reseau} \\
Dans quel ordre les choses se sont-elles passées~? &
Chapitre~\ref{chap:chronologie} \\
Qui a le droit de faire quoi, et selon quelles règles~? &
Chapitre~\ref{chap:normes-droit} \\
Qu'est-ce qu'une vulnérabilité, et comment la coter~? &
Chapitre~\ref{chap:vulnerabilite} \\
Comment les trouver sans attaquer~? & Chapitre~\ref{chap:identifier} \\
Comment décider de ce qu'on corrige d'abord~? & Chapitre~\ref{chap:gerer} \\
Comment remonter de la trace à la faille~? &
Chapitre~\ref{chap:artefact-vuln} \\
Et si l'attaquant a effacé ses traces, ou travaille dans le nuage~? &
Chapitre~\ref{chap:antiforensique} \\
Comment conduire l'enquête de bout en bout~? & Chapitre~\ref{chap:enquete} \\
Comment l'écrire, et la défendre devant un juge~? &
Chapitre~\ref{chap:rapport} \\
Comment mesurer ce que l'enquête ne sait pas~? &
Chapitre~\ref{chap:auditologie} \\
\end{tableaucours}
\vspace{-0.8em}
\begin{center}{\small\itshape Tableau~0.4. Les questions et les chapitres}\end{center}

% ============================================================================
\section{Les métiers de l'affaire}
\label{sec:p00-metiers}
% ============================================================================

Quatre métiers ont travaillé sur l'affaire NYANGA. Le cours les suit
chapitre après chapitre.

\begin{boxmetier}{RSSI}
Mme~NDOUMBE a préservé avant d'analyser, ouvert le registre des actions et
refusé de coter une hypothèse comme un fait. Son rôle, dans les jours qui
suivent, est d'arbitrer entre la reprise de l'activité et la conservation
des preuves, et de faire corriger la faille.
\end{boxmetier}

\begin{boxmetier}{Analyste en réponse à incident}
Mme~MBARGA et son équipe contiennent l'incident, recueillent les pièces,
reconstituent la séquence et remettent à l'établissement un rapport qui
sert à décider. Ils travaillent vite, mais chaque pièce qu'ils touchent
peut finir devant un tribunal.
\end{boxmetier}

\begin{boxmetier}{Application OPJ/Expert Judiciaire}
L'adjudant-chef NANA, officier de police judiciaire, conduit l'enquête
pénale~; M.~ABOUBAKAR, réquisitionné par le procureur, examine les pièces
pour la justice. L'expert ne répond qu'aux questions de sa mission, et
chacune de ses conclusions doit pouvoir être refaite par un autre expert.
\end{boxmetier}

\begin{boxmetier}{Analyste en gestion des vulnérabilités}
Personne n'occupait ce rôle à NYANGA CRÉDIT. C'est la leçon de l'affaire~:
un actif hors de l'inventaire est un actif que personne ne corrige.
\end{boxmetier}

% ============================================================================
\section{Cinq réflexes}
\label{sec:p00-reflexes}
% ============================================================================

\begin{enumerate}
    \item Isoler sans éteindre, puis préserver dans l'ordre de la
          volatilité.
    \item Consigner chaque geste~: heure, auteur, action, résultat,
          empreinte.
    \item Séparer la constatation de l'analyse, et coter \NEK{} ce que les
          pièces ne permettent pas de trancher.
    \item Remonter jusqu'à la faiblesse exploitée~: une enquête qui s'arrête
          aux auteurs laisse la porte ouverte.
    \item Vérifier que la correction est faite, sur la machine elle-même, et
          pas sur la foi d'un courriel.
\end{enumerate}

\separateur

\begin{boxresume}
\textbf{Ce qu'il faut retenir de ce chapitre~:}
\begin{itemize}
    \item Les gestes de la première heure décident des preuves
          disponibles~: isoler, préserver dans l'ordre de la volatilité,
          consigner.
    \item Une constatation se rattache à une source~; une analyse porte un
          niveau de confiance. Une hypothèse reste \NEK{} tant qu'aucune
          pièce ne la tranche.
    \item À NYANGA CRÉDIT, l'enquête remonte du virement frauduleux à une
          faille corrigeable depuis 89~jours, absente de l'inventaire.
    \item L'investigation numérique et l'analyse de vulnérabilités se
          répondent~: l'une explique l'incident, l'autre empêche le
          suivant.
\end{itemize}
\end{boxresume}

\begin{boxinfo}
\textbf{Points de contrôle.} Vous avez compris ce prologue si vous savez~:
\begin{itemize}
    \item citer trois gestes qui préservent la preuve et deux qui la
          détruisent~;
    \item réécrire une phrase de rapport en séparant la constatation et
          l'analyse~;
    \item coter une conclusion \OUI{}, \NON{} ou \NEK{} et justifier la
          cote par une source.
\end{itemize}
\end{boxinfo}

% ============================================================================
\section*{Exercices}
% ============================================================================

\begin{boxexo}[title={\ding{45}\quad Exercice~0.1. Préserver ou détruire \niveaubase}]
Pour chacun des gestes suivants, dites s'il préserve ou s'il détruit de la
preuve, et laquelle.
\begin{enumerate}[(a)]
    \item Lancer un antivirus complet sur le serveur de fichiers.
    \item Photographier l'écran du serveur avant toute action.
    \item Redémarrer le contrôleur de domaine pour « couper l'attaquant ».
    \item Exporter les journaux de la passerelle sur un support neuf et en
          calculer l'empreinte.
    \item Changer immédiatement le mot de passe de \texttt{svc\_backup}
          depuis le serveur compromis.
\end{enumerate}
\end{boxexo}

\begin{boxexo}[title={\ding{45}\quad Exercice~0.2. Constatation ou analyse \niveaumoyen}]
Le premier compte rendu du DSI contient les phrases suivantes. Pour
chacune, dites si c'est une constatation, une analyse ou une hypothèse, et
proposez une cote \OUI{}, \NON{} ou \NEK{} avec sa source.
\begin{enumerate}[(a)]
    \item « Les virements ont été faits depuis l'interface de paiement avec
          notre clé. »
    \item « Un employé a forcément donné son mot de passe. »
    \item « Le serveur de fichiers a envoyé 3,2~Go vers l'extérieur. »
    \item « Les fichiers volés sont ceux des clients. »
    \item « L'attaque vient d'un groupe étranger. »
\end{enumerate}
\end{boxexo}

\begin{boxexo}[title={\ding{45}\quad Exercice~0.3. Le compte des jours \niveaumoyen}]
\begin{enumerate}[(a)]
    \item Vérifiez les deux durées données dans le chapitre~: 77~jours entre
          le correctif et le premier accès, 89~jours entre le correctif et
          l'exfiltration.
    \item La faille a été inscrite au catalogue des vulnérabilités
          exploitées le 19~décembre 2024. Que change cette inscription pour
          un établissement qui suit ce catalogue~?
    \item Donnez deux raisons pour lesquelles un actif peut échapper à
          l'inventaire d'un outil de balayage.
\end{enumerate}
\end{boxexo}

\begin{boxexo}[title={\ding{45}\quad Exercice~0.4. Deux enquêtes, deux obligations \niveauexpert}]
Le cabinet SANKOFA EXPERTISE et l'expert judiciaire M.~ABOUBAKAR examinent
les mêmes pièces.
\begin{enumerate}[(a)]
    \item Pour chacun, dites à qui il rend compte et quelle question il doit
          trancher.
    \item Le directeur général demande au cabinet de ne pas mentionner le
          retard de correction dans son rapport. Que doit faire Mme~MBARGA~?
          L'expert judiciaire serait-il tenu par la même demande~?
    \item Rédigez, pour la cellule de crise, une note de cinq lignes qui
          distingue ce qui est établi, ce qui est réfuté et ce qui reste
          \NEK{} au soir du 14~mars.
\end{enumerate}
\end{boxexo}

\bigskip

\begin{boxinfo}
\textbf{Pour aller plus loin.}
\begin{itemize}
    \item La RFC~3227~\cite{rfc3227}, pour l'ordre de volatilité.
    \item Le NIST SP~800-86~\cite{nist80086}, pour l'intégration de la
          forensique à la réponse aux incidents.
    \item Le catalogue KEV de la CISA~\cite{cisakev}.
\end{itemize}
\end{boxinfo}

\begin{boxinfo}
\textbf{Transition.} Avant de toucher aux pièces de l'affaire, il faut
savoir ce qu'est une enquête numérique~: son objet, ses principes, ses
étapes. C'est l'objet du Chapitre~\ref{chap:investigation}.
\end{boxinfo}
